\section{Limiti e sviluppi futuri}
\label{sec:limiti}

\subsection{Limiti dell'implementazione attuale}
\label{sec:limiti-attuali}

Il modello è verificato sul bagno termico. Per usarlo in un caso con flusso
mancano ancora alcuni pezzi, e restano alcune differenze rispetto al lavoro di
riferimento:

\begin{itemize}[leftmargin=1.6em]
  \item \textbf{Trasporto dell'energia vibro-elettronica.} La sua equazione
    non contiene ancora il trasporto con il flusso né la conduzione del
    calore, e il campo non ha condizioni al contorno: nel bagno termico non
    servono, in un caso con flusso sì. Allo stesso modo, le sorgenti usano la
    densità del passo precedente invece di quella appena aggiornata, cosa
    indifferente finché il gas è fermo.
  \item \textbf{Proprietà di trasporto.} Viscosità e conducibilità vengono
    ancora dal modello a una temperatura del template; nel bagno termico non
    intervengono, ma in un flusso viscoso andrebbero calcolate con il modello
    a due temperature.
  \item \textbf{Modello a polinomi del template.} Resta attivo solo per
    costruire i campi, ma sopra i 20\,000~K emette avvisi. Sono innocui,
    perché temperature e densità vengono da \mpp, ma andrebbero eliminati.
  \item \textbf{Differenze di modello.} Nel solver c'è una sola temperatura
    vibrazionale, e lo scambio V--V esiste solo nei programmi
    zero-dimensionali; i tempi di rilassamento sono mediati come in \mpp;
    l'accoppiamento fra chimica e vibrazione è solo quello non preferenziale;
    il procedimento per l'esponente di Park vale solo se il meccanismo
    contiene soltanto dissociazioni.
  \item \textbf{Costo di calcolo.} Sullo stesso caso il solver è oltre mille
    volte più lento dei programmi zero-dimensionali. Parte del costo viene dal
    fatto che le velocità di reazione sono ricalcolate più volte per cella a
    ogni passo, separatamente per ogni specie e per ogni termine sorgente che
    le usa.
\end{itemize}

\subsection{Sviluppi futuri}
\label{sec:sviluppi}

\begin{enumerate}[leftmargin=1.8em]
  \item Completare l'equazione dell'energia vibro-elettronica con il
    trasporto, la conduzione e le condizioni al contorno, e passare a un caso
    con flusso. Il primo passo naturale è il tubo d'urto monodimensionale, per
    il quale il template contiene già un caso di partenza.
  \item Calcolare viscosità e conducibilità con \mpp, così da poter trattare
    flussi viscosi.
  \item Affrontare i casi bi- e tridimensionali della seconda parte del lavoro
    di riferimento.
  \item Ridurre le differenze di modello: implementare la media dei tempi di
    rilassamento del lavoro di riferimento, per verificare se il ritardo della
    temperatura vibrazionale nelle miscele scompare, e aggiungere più
    temperature vibrazionali e modelli di accoppiamento più raffinati.
  \item Ridurre il costo, calcolando le velocità di reazione una sola volta
    per cella e per passo.
\end{enumerate}
